% ---------------------------------------------------------------
\usepackage[margin=2.4cm]{geometry}

% ---------------------------------------------------------------
% ENCODING AND FONTS
% ---------------------------------------------------------------
\usepackage[utf8]{inputenc}
\usepackage[T1]{fontenc}
\usepackage{lmodern}           % CM-compatible outlined fonts
% Alternative: Times-like math (closer to AMS / IEEE style)
%   \usepackage{newtxtext,newtxmath}  % requires newtx package
\usepackage{microtype}         % character-level justification; always last font pkg
\linespread{1.04}              % slight openness; 1.0 is pure single-spacing

% ---------------------------------------------------------------
% MATHEMATICS — core
% ---------------------------------------------------------------
\usepackage{amsmath,amssymb,amsthm,mathtools}
\usepackage{bm}                % bold math (vectors, matrices)
\usepackage{physics}           % \abs, \norm, \dd, \grad, …
\usepackage{mathrsfs}          % \mathscr{} script letters (e.g. for operators)
\usepackage{bbm}               % \mathbbm{1} blackboard-bold indicator
\usepackage{cancel}            % \cancel, \bcancel for struck-through terms
\usepackage{xfrac}             % \sfrac{}{} for inline slanted fractions

% Allow long equation arrays to break across pages
\allowdisplaybreaks[2]

% ---------------------------------------------------------------
% MATHEMATICS — highlighted / boxed equations
% ---------------------------------------------------------------
\usepackage{empheq}            % \begin{empheq}[box=\fbox]{align}
% Define a shaded box option for main results:
\usepackage[most]{tcolorbox}
\tcbset{highlight math style={%
    enhanced, colframe=black!70, colback=gray!6,
    arc=2pt, boxrule=0.5pt, left=3pt, right=3pt}}

% ---------------------------------------------------------------
% SI UNITS
% ---------------------------------------------------------------
\usepackage{siunitx}           % \SI{28}{\giga\hertz}, \si{\watt\per\kilo\gram}
\sisetup{
  detect-all,                  % inherit surrounding font
  range-phrase = {--},
  range-units  = single,
  per-mode     = symbol,
  inter-unit-product = {\,}
}

% ---------------------------------------------------------------
% FIGURES AND TABLES
% ---------------------------------------------------------------
\usepackage{graphicx}
\graphicspath{{figures/}}
\usepackage{booktabs}          % \toprule, \midrule, \bottomrule
\usepackage{array}             % extended column specifiers
\usepackage{tabularx}          % full-width tables
\usepackage{multirow}
\usepackage{makecell}          % multi-line table cells
\usepackage{caption}
\captionsetup{font=small, labelfont=bf, labelsep=period,
              skip=4pt}        % compact captions, bold label, period separator
\usepackage{subcaption}
\usepackage{threeparttable}    % table notes for lit-review tables
\usepackage{float}
\usepackage[section]{placeins} % keep floats within their section

% Relax float placement rules
\renewcommand{\topfraction}{0.9}
\renewcommand{\bottomfraction}{0.9}
\renewcommand{\textfraction}{0.1}
\renewcommand{\floatpagefraction}{0.7}
\setcounter{topnumber}{4}
\setcounter{bottomnumber}{4}
\setcounter{totalnumber}{8}

% ---------------------------------------------------------------
% SECTION HEADINGS (titlesec)
% ---------------------------------------------------------------
\usepackage{titlesec}
\titleformat{\section}{\large\bfseries}{\thesection}{1em}{}
\titleformat{\subsection}{\normalsize\bfseries}{\thesubsection}{1em}{}
\titleformat{\subsubsection}{\normalsize\bfseries}{\thesubsubsection}{1em}{}
\titlespacing*{\section}{0pt}{1.6ex plus .4ex minus .2ex}{0.8ex plus .2ex}
\titlespacing*{\subsection}{0pt}{1.2ex plus .3ex minus .2ex}{0.5ex plus .1ex}

% ---------------------------------------------------------------
% MISCELLANEOUS
% ---------------------------------------------------------------
\usepackage{enumitem}          % fine-grained list control
\usepackage{xcolor}
\usepackage{xspace}            % smart spacing after macro names
\usepackage{todonotes}[disable]% \todo{} comments — set [disable] for final version
\usepackage{fancyhdr}          % running headers

% Running header: empty header, page number centred in footer
\pagestyle{fancy}
\fancyhf{}
\cfoot{\thepage}
\renewcommand{\headrulewidth}{0pt}   % no rule when header is empty

% ---------------------------------------------------------------
% BIBLIOGRAPHY
% ---------------------------------------------------------------
\usepackage[numbers,sort&compress]{natbib}

% ---------------------------------------------------------------
% HYPERLINKS — load near last, before cleveref
% ---------------------------------------------------------------
\usepackage[
  colorlinks = true,
  linkcolor  = {blue!55!black},
  citecolor  = {green!45!black},
  urlcolor   = {blue!65!black},
  pdfstartview = FitH,
  bookmarksnumbered = true,
  bookmarksopen     = true
]{hyperref}
\usepackage{bookmark}          % fixes bookmark levels in PDF outline
\usepackage{doi}               % live DOI hyperlinks (must load AFTER hyperref)

% cleveref must come AFTER hyperref
\usepackage[nameinlink,capitalize]{cleveref}

% ---------------------------------------------------------------
% THEOREM-LIKE ENVIRONMENTS
% (all share the same counter as theorem, numbered per section)
% ---------------------------------------------------------------
\usepackage{aliascnt}

\theoremstyle{definition}      % upright body text for all env.

\newtheorem{theorem}{Theorem}[section]

\newaliascnt{proposition}{theorem}
\newtheorem{proposition}[proposition]{Proposition}
\aliascntresetthe{proposition}

\newaliascnt{lemma}{theorem}
\newtheorem{lemma}[lemma]{Lemma}
\aliascntresetthe{lemma}

\newaliascnt{corollary}{theorem}
\newtheorem{corollary}[corollary]{Corollary}
\aliascntresetthe{corollary}

\newaliascnt{definition}{theorem}
\newtheorem{definition}[definition]{Definition}
\aliascntresetthe{definition}

\newaliascnt{remark}{theorem}
\newtheorem{remark}[remark]{Remark}
\aliascntresetthe{remark}

% cleveref names for theorem-like environments
\crefname{proposition}{Proposition}{Propositions}
\Crefname{proposition}{Proposition}{Propositions}
\crefname{lemma}{Lemma}{Lemmas}
\Crefname{lemma}{Lemma}{Lemmas}
\crefname{corollary}{Corollary}{Corollaries}
\Crefname{corollary}{Corollary}{Corollaries}
\crefname{definition}{Definition}{Definitions}
\Crefname{definition}{Definition}{Definitions}
\crefname{remark}{Remark}{Remarks}
\Crefname{remark}{Remark}{Remarks}

% ---------------------------------------------------------------
% CUSTOM MATH COMMANDS
% ---------------------------------------------------------------
% --- vectors and fields ---
\newcommand{\khat}{\hat{\bm{k}}}
\newcommand{\nhat}{\hat{\bm{n}}}
\newcommand{\rr}{\mathbf{r}}
\newcommand{\EE}{\mathbf{E}}
\newcommand{\HH}{\mathbf{H}}
% --- dosimetric quantities ---
\newcommand{\Sinc}{S_{\mathrm{inc}}}
\newcommand{\Sab}{S_{\mathrm{ab}}}
\newcommand{\Teff}{T_{\mathrm{eff}}}
\newcommand{\Tavg}{T_{\mathrm{avg}}}
\newcommand{\Aab}{A_{\mathrm{ab}}}
\newcommand{\Aperp}{A_{\perp}}
\newcommand{\SARwb}{\mathrm{SAR}_{\mathrm{wb}}}
% --- material parameters ---
\newcommand{\ntilde}{\tilde{n}}
\newcommand{\epstilde}{\tilde{\varepsilon}_r}
% --- operators ---
\DeclareMathOperator{\ReLU}{ReLU}
% (\erf already provided by the physics package)
\DeclareMathOperator{\RE}{Re}
\DeclareMathOperator{\IM}{Im}
% --- other shorthands ---
\newcommand{\Tbar}{\bar{T}}
\newcommand{\diff}{\mathrm{d}}        % upright differential
\newcommand{\ident}{\mathbbm{1}}       % indicator / identity  (requires bbm)
\newcommand{\Reals}{\mathbb{R}}
\newcommand{\Complex}{\mathbb{C}}
\newcommand{\half}{\tfrac{1}{2}}

% ---------------------------------------------------------------
% PART HEADINGS (article class has no \part)
% ---------------------------------------------------------------
